% =====================================================================
%  作品简介（中文参考版）—— 格致
%  ⚠️ 仅供内部评审与指导老师审阅，切勿提交。
%     竞赛章程规定所有提交材料只接受英文，正式提交请用
%     ../ProjectIntroduction_Gezhi.pdf
%
%  编译：xelatex ProjectIntroduction_Gezhi_中文.tex   （跑两遍）
% =====================================================================
\documentclass[11pt,a4paper,UTF8,fontset=windows]{ctexart}

\usepackage[margin=2.0cm,top=1.9cm,bottom=1.9cm]{geometry}
\usepackage{booktabs}
\usepackage{tabularx}
\usepackage{array}
\usepackage{enumitem}
\usepackage{titlesec}
\usepackage{xcolor}
\usepackage{caption}
\usepackage{hyperref}
\usepackage{xurl}

\newcolumntype{R}{>{\raggedright\arraybackslash}X}
\newcommand{\fillin}[1]{\textcolor{red}{\textbf{【待填：#1】}}}
\newcommand{\fillc}[1]{\textcolor{red}{\textbf{【#1】}}}

\titleformat{\section}{\normalfont\large\bfseries}{\thesection.}{0.5em}{}
\titlespacing*{\section}{0pt}{8pt}{4pt}
\titleformat{\subsection}[runin]{\normalfont\bfseries}{}{0em}{}[\hspace{0.5em}]
\titlespacing*{\subsection}{0pt}{5pt}{0.5em}

\setlist[itemize]{leftmargin=2em,itemsep=1pt,topsep=2pt,parsep=0pt}
\setlist[enumerate]{leftmargin=2.4em,itemsep=1pt,topsep=2pt,parsep=0pt}
\setlength{\parindent}{2\ccwd}
\setlength{\parskip}{0pt}
\linespread{1.22}
\captionsetup{font=small,labelfont=bf}
\pagestyle{empty}

\begin{document}

% ===================== 标题 =========================================
\begin{center}
{\LARGE\bfseries 格致：一个把学生放上「教师」位置的\\[2pt]
本地优先 AI 学习伙伴}\\[7pt]
{\small 2026 年全球青少年人工智能未来创新竞赛 · 澳门中学生赛区\\
主题：AI 与教育}\\[4pt]
{\footnotesize 队伍 \textit{[队伍名称]} $\cdot$ \textit{[学校名称]} $\cdot$ 澳门}
\end{center}
\vspace{-2pt}
\hrule
\vspace{5pt}

% ===================== 一、问题 =====================================
\section{问题}

生成式人工智能被寄予厚望，被认为终于能让每个学生都拥有一位私人导师。但最早的一批
严谨证据给出的答案要复杂得多，而且两个方向的结论同时成立。

\begin{table}[h]
\centering
\small
\begin{tabularx}{\textwidth}{@{}p{2.6cm} p{3.5cm} R@{}}
\toprule
\textbf{研究} & \textbf{设计} & \textbf{结论} \\
\midrule
Bastani 等（2024） &
随机对照实验；土耳其约 1000 名高中生；数学 &
可无限制使用 ChatGPT 的组：练习正确率 \textbf{+48\%}，但撤掉 AI 后考试
\textbf{低 17\%}。使用"只给提示、不给答案"的导师版：练习 \textbf{+127\%}，
但考试\textbf{与不用 AI 的对照组没有差别}。 \\[2pt]
Oreopoulos 与 Low（2026） &
两年整群随机对照实验；美国田纳西州 18 所初中；AI 被配置成只引导不给答案 &
效应仅 \textbf{0.06--0.08 个标准差}，与"同一平台不用 AI"相当。瓶颈是行为而非技术：
中位数学生只在\textbf{三分之一的练习日}使用它，在做错题的场次中只有
\textbf{17\%} 会去求助。 \\
\bottomrule
\end{tabularx}
\end{table}

\vspace{-4pt}
把这两项研究放在一起看，指向的是同一条窄缝：\textbf{今天的 AI 导师只有两种结局——
要么有害，要么无人使用。} 而且"把 AI 做得更克制"并不能解决第二个问题：一个被刻意
设计成只给提示的导师同样没有产生可测量的收益，因为学生根本不去打开它。唯一从未被
证明会损害学习的条件，是\textbf{由学生自己产出解释}的那一个。

% ===================== 二、目标使用者 ===============================
\section{目标使用者}

\textbf{主要使用者：}澳门的中学生（中三至中六，约 14--18 岁），正在准备校内考试与
大学入学。我们具体瞄准的是这样一种学生——他\emph{感觉}自己懂了，却讲不出来。这种
"懂了错觉"并非个例：综合 115 项研究、15889 名学习者、502 个效应量，学生对自己掌握
程度的判断与实际成绩之间的加权平均相关只有 \(r = .178\)（Yang 等，2022）。

\textbf{次要使用者：}需要知道一个班\emph{真正}理解了什么（而不只是做完了什么）的学科
教师；以及希望看到比"学习时长"更有信息量的东西的家长。

% ===================== 三、主要功能 =================================
\section{主要功能}

格致是一款 Windows 桌面学习应用，整款产品围绕一个循环组织：
\textbf{先暴露无知，再产出解释，被质疑，然后掌握。}

\begin{itemize}
  \item \textbf{知识库。}学生自己的课程材料——教材、讲义、笔记、扫描的练习卷——
  在本地建立索引，中英文皆可检索。任何内容都不上传。
  \item \textbf{题库与组卷。}题目按知识点标注，可导出为排版规范的试卷；另有
  \emph{交错模式}，把不同知识点混排，并刻意\emph{不标注}每道题属于哪个知识点。
  \item \textbf{讲述与答辩（核心）。}学生把某个知识点讲给屏幕上的伙伴「小格」听。
  小格被刻意设定成\emph{比他低一个年级}，并且带着一套固定的错误理解。它不会点头
  附和，而是会反驳。系统按六维评分表打分，并记录"用了几个回合才真正说服它"。
  \item \textbf{间隔复习与计划。}每个知识点带一个掌握度估计，并据此同时排程
  \emph{题目}和\emph{讲解}。
  \item \textbf{不看分数也能看到进步。}小格头顶有一个九宫格。九格全空表示它什么都
  还不懂；每成功讲清一次，就亮起一格。学生不必读数字就知道自己走到哪里了。
\end{itemize}

% ===================== 四、AI 方法 ==================================
\section{AI 方法}

格致没有自己实现 Agent 循环。推理层采用 \textbf{DeepSeek Harness（DSH）}——一个
MIT 许可的 Agent 内核，以受控子进程方式调用，并输出按行分隔的 JSON 事件流。会话
持久化、工具沙箱、上下文压缩与审批策略由此直接复用，这些都不是一个中学生团队应当
从零重造的部分。其中三个技术决定最关键：

\textbf{一、知识库本身就是沙箱。}子进程的工作目录被设为学生的知识库，内核据此推导
出写入边界。AI 能读写学生自己的笔记，除此之外什么也碰不到——这条边界由操作系统的
进程模型强制，而不是靠提示词里的嘱咐。

\textbf{二、中文检索需要一个不那么显然的修复。}我们实测发现，SQLite 默认的 FTS5
分词器对中文子串查询返回\emph{零}结果，因为它把一整句中文当成一个词元。改用
\texttt{tokenize='trigram'} 后恢复正常：在 200 份文档上查询延迟约 1\,ms。因此检索
完全在本地、离线进行，并且是可溯源的——系统给出的每一条论断都锚定在学生自己教材的
某一页上。

\textbf{三、这个伙伴是被刻意设计成会犯错的。}一份误解库（来自学生自己的错题统计）
被注入小格的人格提示词，同时加了一条硬约束：\textbf{在学生的解释真正命中该误解的
纠正要点之前，它不得表示理解。}这一设计依据两项研究：CHI 2025 的工作表明可教代理
必须\emph{装笨}才有教学价值；而另一项对照实验发现，当 AI 直接给出正确答案时，学生
的纠错能力\textbf{完全没有提升}。

% ===================== 五、创新点 ===================================
\section{创新点}

我们不声称发明了"以教促学"。可教代理的研究最早可追溯到 Chase 等（2009），2024--
2026 年更是一个活跃方向。我们的贡献是四件具体的事：

\begin{enumerate}
  \item 这类研究绝大多数以\textbf{大学生}为对象，且多在计算机课程中开展。我们把它
  带到了\textbf{中学生}面前，做成一款面向他们自己课程的中文桌面产品。
  \item 学生教的是\textbf{他自己教材里的知识点}，由本地检索提供并附页码引用——
  而不是通用题目。
  \item 我们把"说服小格"这件事变成了\textbf{可量化、并能喂给间隔复习排程的分数}，
  形成闭环，而不是一次性的聊天。
  \item 知识库与学习记录\textbf{从不离开本机}；模型调用走用户自己的账号，提问内容
  不做留存。
\end{enumerate}

% ===================== 六、预期影响 =================================
\section{预期影响}

我们以学习结果而非使用时长来衡量成败，用三个指标：(1) 一个知识点在"讲述—答辩"前后
费曼评分表总分的变化；(2) 学生自评分与评分表得分之间的\textbf{校准偏差}，并做长期
追踪；(3) 消除一个误解所需的对话回合数——理解越深，回合数应越低。

格致回应的是\textbf{可持续发展目标 4（优质教育）}：把通常只有一对一补习才能获得的
"解释性练习"带给任何一个有笔记本电脑的学生；同时把\textbf{AI 素养}本身当作一个学习
目标——能诊断并纠正 AI 错误理解的学生，才算真正知道这类系统会在哪里出错。这一点
直接对应 UNESCO《学生人工智能能力框架》（2024）中"AI 系统设计"维度的"创造"层级。

% ===================== 七、目前进展 =================================
\section{目前进展}

一个可运行的原型已在 Windows 上跑通上述完整闭环：中英混合材料的本地索引、带引用的
检索、与小格的语音讲述会话、评分表打分，以及导出排版试卷。误解库目前有
\fillin{条数} 条，覆盖 \fillin{科目}。我们已经完成第一轮 \fillin{人数}
名中学生的试用，实测结果见配套的《研究报告》。

下一步是加入"不使用 AI"的对照组、用真人评分者校验评分表，并把研究周期拉长到间隔
复习排程真正能发挥作用的时间尺度。

\end{document}
